\documentclass[11pt]{article}
\usepackage{amsmath,amssymb,amsthm,mathtools}
\usepackage[margin=1in]{geometry}
\usepackage{booktabs,longtable,array}
\usepackage{hyperref}
\setlength{\emergencystretch}{3em}

\newtheorem{definition}{Definition}
\newtheorem{assumption}{Assumption}
\newtheorem{theorem}{Theorem}
\newtheorem{lemma}{Lemma}
\newtheorem{requiredlemma}{Required Lemma}
\newtheorem{proposition}{Proposition}
\newtheorem{corollary}{Corollary}
\newtheorem{remark}{Remark}

\title{Technical Gates for Amplification Payment:\\
K3--K7 Source-Ledger Burden}
\author{Joseph Roy}
\date{June 2026}

\begin{document}
\maketitle

\begin{abstract}
This note isolates the mathematical proof burden that is imported by the public
paper \emph{Amplification Must Be Paid}.  It is a trimmed technical ledger for
the downstream gates K3--K7: quantitative quotient-or-charge, forward
absorption, reverse matrix absorption, high-strain transport rigidity, spectral
and loop-gap certification, and compatibility with the regularity bridge.

The note starts after the upstream source-promotion data have been granted:
there is a scale sequence \(r_n\downarrow0\), a candidate point \(z_0\), a
positive source ledger \(P(r_n)\), and promoted finite-thickness source
families with fixed capture fraction \(\eta>0\).  The purpose here is to state
what K3--K7 must prove, how the constants feed the loop, and where each open
analytic burden remains.
\end{abstract}

\clearpage
\section{Scope and Inputs}
\label{sec:scope}

This file is deliberately narrower than the full source-packing program.  It
does not re-prove source persistence, source promotion, or construction of the
promoted family.  It assumes that those upstream objects have been supplied and
then asks whether the downstream ledger can close.

\begin{assumption}[Upstream promoted-source package]
\label{ass:upstream-package}
Along a scale sequence \(r_n\downarrow0\) at a candidate point \(z_0=(x_0,T)\),
assume:
\begin{enumerate}
\item \(P(r_n)\ge p_*>0\);
\item a promoted finite-thickness source family captures a fixed fraction
\(\eta>0\) of the source;
\item the promoted family has the usual core/collar/tail decomposition and
finite-overlap bookkeeping;
\item quotient redundancy, if it occurs, is a genuine identification of the
same physical source and not a new nonredundant amplification channel.
\end{enumerate}
\end{assumption}

The downstream proof burden is:
\[
\text{K3}+\text{K4}+\text{K5/K5B}+\text{K6}+\text{K7}.
\]
The constant ledger at the end of the note shows why qualitative positivity is
not enough: the loop closes only under a strict numerical gap.

\clearpage
\section{Definitions and Denominator}
\label{sec:definitions}

Fix
\[
Q_r(z_0):=B_r(x_0)\times(T-r^2,T),
\qquad
\sigma_+:=(\omega\cdot S\omega)_+,
\]
and define the source ledger
\begin{equation}
\label{eq:P-def}
P(r):=r\iint_{Q_r(z_0)}\sigma_+\,dx\,dt .
\end{equation}

\begin{assumption}[Local source-ledger finiteness]
\label{ass:source-finite}
For every fixed parent cylinder used in the ledger,
\[
P(r_0)=r_0\iint_{Q_{r_0}(z_0)}\sigma_+\,dx\,dt<\infty.
\]
This is automatic in smooth pre-singular cylinders.  In a suitable-weak
formulation it is a standing local source-finiteness or truncation hypothesis.
\end{assumption}

The scale-critical strain variable is
\[
S_r(x,t):=r^2|S_{\mathrm{loc}}(x,t)|,
\]
and the strain row is
\begin{equation}
\label{eq:strain-row}
\mathfrak B_{\mathrm{strain}}(r)
:=
r\iint_{\mathcal U_r}\sigma_+ S_r\phi^2\,dx\,dt
=
r^3\iint_{\mathcal U_r}\sigma_+|S_{\mathrm{loc}}|\phi^2\,dx\,dt .
\end{equation}
The unnormalized term
\[
r\iint\sigma_+|S_{\mathrm{loc}}|\phi^2
\]
is supercritical and is not part of the live denominator.

Use the row vector
\begin{equation}
\label{eq:X-def}
X(r):=
\begin{pmatrix}
\mathcal E_\omega\\
\mathcal B_\omega\\
\mathcal D_{\mathrm{ARR}}\\
\mathfrak B_{\mathrm{strain}}\\
M_{\mathrm{deact}}\\
\mathcal E_\nu^{\mathrm{coh},+}
\end{pmatrix},
\qquad
D(r):=\mathbf 1^T X(r).
\end{equation}
The rows are vorticity energy-dissipation, boundary/collar flux, annular
renewal/tail reservoir, source-weighted strain, deactivation, and positive
coherent viscous residual.

\begin{definition}[Packed boundary row]
\label{def:packed-boundary}
The boundary row \(\mathcal B_\omega(R;z_0)\) is parent-packed, not merely a
single shell.  It contains a one-scale shell row and the logarithmic collar
packing
\[
\mathcal B_{\omega,\lambda_0}^{\mathrm{pack}}(R;z_0)
:=
\sup_{\substack{0<a<b\\ \lambda_0b\le R}}
\int_a^b
\mathcal B_\omega^{\mathrm{adv}}(s,\lambda_0;z_0)\frac{ds}{s},
\]
where
\[
\mathcal B_\omega^{\mathrm{adv}}(s,\lambda_0;z_0)
:=
s\iint_{Q_{\lambda_0s}(z_0)\setminus Q_s(z_0)}
|u-\bar u_{B_{\lambda_0s}}(t)|\,|\omega|^2\,
|\nabla\phi_{s,\lambda_0}|\,\phi_{s,\lambda_0}\,dx\,dt .
\]
This packed convention is part of K5 and K7 compatibility.
\end{definition}

The mean-subtracted CKN quantity is
\begin{equation}
\label{eq:Kms}
\mathcal K_{\mathrm{CKN}}^{\mathrm{ms}}(r;z_0)
:=
r^{-2}\iint_{Q_r(z_0)}
\Big(|u-\bar u_{B_r}(t)|^3
+|p-\bar p_{B_r}(t)|^{3/2}\Big)\,dx\,dt .
\end{equation}

\clearpage
\section{K3: Quantitative Quotient-or-Charge}
\label{sec:K3}

The K3 burden is to turn nonredundant promoted packing into a positive final
physical charge.  The final physical ledger must include oscillation/twist,
because large accumulated strain need not become net shape growth.

On each active branch core choose a source director \(e_j\in S^2\) and use the
sign-free projector
\[
E_j:=e_j\otimes e_j .
\]
Define the branch oscillation density
\begin{equation}
\label{eq:theta-osc}
\Theta_{\mathrm{osc}}^{(j)}
:=
\min\{1,\ r^2|D_tE_j|+r|\nabla E_j|\},
\qquad
D_t:=\partial_t+u\cdot\nabla .
\end{equation}
On an overlap \(\Omega_{ij}:=\mathcal T_i\cap\mathcal T_j\cap Q_r\), set
\begin{equation}
\label{eq:Phi-osc}
\Phi_{\mathrm{osc}}(i,j)
:=
\min\Bigg\{1,\,
\frac{
r\iint_{\Omega_{ij}}\sigma_+
\big(|E_i-E_j|+\Theta_{\mathrm{osc}}^{(i)}
+\Theta_{\mathrm{osc}}^{(j)}\big)\phi^2\,dx\,dt
}{
m_{ij}+\varepsilon
}\Bigg\},
\end{equation}
where
\[
m_{ij}:=r\iint_{\Omega_{ij}}\sigma_+\phi^2\,dx\,dt.
\]
The family oscillation ledger is the source-weighted average
\begin{equation}
\label{eq:Rosc-family}
R_{\mathrm{osc}}^{\mathrm{fam}}(r)P(r)
:=
\sum_{j}P_j(r)R_{\mathrm{osc}}^{(j)}(r)+o_r(1).
\end{equation}
The final physical cost is
\begin{equation}
\label{eq:Fphys-final}
\mathcal F_{\mathrm{phys}}^{\mathrm{final}}
:=
\mathcal F_{\mathrm{phys}}+R_{\mathrm{osc}}^{\mathrm{fam}}.
\end{equation}

\begin{definition}[Quantitative final quotient redundancy]
\label{def:quant-final-redundancy}
Fix a quotient threshold \(\delta>0\) and let
\(\sim_\delta^{\mathrm{final}}\) be the equivalence relation generated by
significant-overlap pairs with
\(d_{\mathrm{phys}}^{\mathrm{final}}\le\delta\), where the final edge distance
is defined in \eqref{eq:dphys-final} below.  A promoted family is
\((\delta,\eta)\)-quotient-redundant at scale \(r\) if, after finite-overlap
normalization, at least \(\eta P(r)\) of the captured promoted source lies
inside these quotient classes and the nonclass edge mass satisfies
\[
\sum_{(i,j):\,i\not\sim_\delta^{\mathrm{final}}j}m_{ij}
<c_m^{\mathrm{final}}\eta P(r).
\]
Thus redundancy is not a verbal escape clause: it is the branch where captured
mass is quantitatively accounted for by physical identification.
\end{definition}

\begin{proposition}[Oscillation ledger stabilization]
\label{prop:osc-stabilization}
\(R_{\mathrm{osc}}^{\mathrm{fam}}\) is a separate physical coordinate, not a
shape, renewal, spread, or activity proxy.  Moreover, \(\Phi_{\mathrm{osc}}\)
is source-adjoint to \(R_{\mathrm{osc}}^{\mathrm{fam}}\): a heavy
nonredundant overlap with \(\Phi_{\mathrm{osc}}\ge\delta\) contributes a
scale-independent positive fraction of oscillation density to the family
ledger, up to finite-overlap and quotient-chain constants.
\end{proposition}

\begin{proof}
The old shape ledger is Cauchy--Green/eigenvalue based; spread records support
separation; renewal records branch identity turnover; and activity records
source persistence.  None of these sees \(D_tE_j\), \(\nabla E_j\), or the
projective direction mismatch \(|E_i-E_j|\).  Thus a branch may keep support,
ancestry, activity, and net shape spectrum while rotating its source direction.
Equation~\eqref{eq:Phi-osc} integrates exactly the source-weighted density used
in \eqref{eq:Rosc-family}; finite-overlap normalization converts a heavy
overlap into a family-level lower bound.
\end{proof}

\begin{requiredlemma}[Gate K3-A: finite-overlap quotient dichotomy]
\label{lem:K3-dichotomy}
For promoted families obeying the finite-overlap bookkeeping, exactly one of
the following alternatives must be available with constants independent of
\(r\):
\[
\sum_{(i,j):\,i\not\sim_\delta^{\mathrm{final}}j}m_{ij}
\ge c_m^{\mathrm{final}}\eta P(r),
\]
or the family is \((\delta,\eta)\)-quotient-redundant in the sense of
Definition~\ref{def:quant-final-redundancy}.
\end{requiredlemma}

\begin{proof}[Required route]
Build the finite weighted overlap graph whose vertices are promoted packets and
whose edge weights are \(m_{ij}\).  Collapse the graph along
\(\sim_\delta^{\mathrm{final}}\).  Finite overlap makes total captured mass
comparable to \(\eta P(r)\) up to uniform constants.  If the quotient-collapsed
graph has too little edge mass outside its classes, the captured mass is
concentrated inside quotient classes and is, by definition, the redundant
branch.  The contrapositive gives the displayed nonredundant mass lower bound.
\end{proof}

\begin{requiredlemma}[Gate K3-B: final quantitative quotient-or-charge]
\label{lem:K3}
Define
\begin{equation}
\label{eq:dphys-final}
d_{\mathrm{phys}}^{\mathrm{final}}(i,j)
:=
\max\{d_{\mathrm{phys}}(i,j),\,w_{\mathrm{osc}}\Phi_{\mathrm{osc}}(i,j)\}.
\end{equation}
Let \(\sim_\delta^{\mathrm{final}}\) be generated by significant-overlap pairs
with \(d_{\mathrm{phys}}^{\mathrm{final}}\le\delta\), and set
\[
E_{\mathrm{nonred}}^{\mathrm{final}}
:=\{(i,j):i\not\sim_\delta^{\mathrm{final}}j,\ m_{ij}>0\}.
\]
There are constants \(c_m^{\mathrm{final}}>0\), \(C_*^{\mathrm{final}}>0\),
and \(\delta>0\) such that, if the promoted family is not quotient-redundant
and \(R_{\mathrm{pack},\delta}^{*,\mathrm{final}}\ge\eta\), then
\begin{align}
\sum_{(i,j)\in E_{\mathrm{nonred}}^{\mathrm{final}}}m_{ij}
&\ge c_m^{\mathrm{final}}\eta P(r), \label{eq:K3-mass}\\
\sum_{(i,j)\in E_{\mathrm{nonred}}^{\mathrm{final}}}
m_{ij}d_{\mathrm{phys}}^{\mathrm{final}}(i,j)
&\le
C_*^{\mathrm{final}}P(r)\mathcal F_{\mathrm{phys}}^{\mathrm{final}}(r).
\label{eq:K3-ledger}
\end{align}
Consequently,
\begin{equation}
\label{eq:cfinal}
\mathcal F_{\mathrm{phys}}^{\mathrm{final}}(r)
\ge c_{\mathrm{final}}(\eta),
\qquad
c_{\mathrm{final}}(\eta)
:=
\frac{\delta c_m^{\mathrm{final}}}{C_*^{\mathrm{final}}}\eta,
\end{equation}
or the family is quantitatively quotient-redundant.
\end{requiredlemma}

\begin{proof}[Required route]
Every edge that survives the generated quotient satisfies
\(d_{\mathrm{phys}}^{\mathrm{final}}\ge\delta\).  If the mass lower bound
\eqref{eq:K3-mass} failed, Lemma~\ref{lem:K3-dichotomy} routes the family into
the quantitative quotient-redundant branch.  The old physical coordinates give
the old ledger-adjoint edge estimate, and
Proposition~\ref{prop:osc-stabilization} adds the oscillation estimate
\[
\sum m_{ij}\Phi_{\mathrm{osc}}(i,j)
\le C_{\mathrm{osc},*}P(r)R_{\mathrm{osc}}^{\mathrm{fam}}(r)+o_r(1).
\]
Since \(d_{\mathrm{phys}}^{\mathrm{final}}\le
d_{\mathrm{phys}}+w_{\mathrm{osc}}\Phi_{\mathrm{osc}}\), the edge ledger
\eqref{eq:K3-ledger} follows after enlarging \(C_*^{\mathrm{final}}\).  The
threshold \(d_{\mathrm{phys}}^{\mathrm{final}}\ge\delta\), the mass bound, and
\eqref{eq:K3-ledger} imply
\[
\delta c_m^{\mathrm{final}}\eta P(r)
\le C_*^{\mathrm{final}}P(r)\mathcal F_{\mathrm{phys}}^{\mathrm{final}}(r),
\]
and cancellation on the nontrivial source sequence gives \eqref{eq:cfinal}.
\end{proof}

\clearpage
\section{K4: Forward Absorption}
\label{sec:K4}

K4 routes the final physical charge into the denominator \(D\).

\begin{lemma}[Old forward source absorption]
\label{lem:K4-old}
There is \(C_{\mathrm{old}}<\infty\) such that
\begin{equation}
\label{eq:K4-old}
P(r)\mathcal F_{\mathrm{phys}}(r)
\le C_{\mathrm{old}}D(r)+o_r(1).
\end{equation}
\end{lemma}

\begin{lemma}[Oscillation/twist forward absorption]
\label{lem:K4-osc}
There is \(C_{\mathrm{osc}}<\infty\) such that
\begin{equation}
\label{eq:K4-osc}
P(r)R_{\mathrm{osc}}^{\mathrm{fam}}(r)
\le C_{\mathrm{osc}}D(r)+o_r(1).
\end{equation}
\end{lemma}

\begin{proof}
Using \eqref{eq:theta-osc} and \eqref{eq:Rosc-family},
\[
P R_{\mathrm{osc}}^{\mathrm{fam}}
\lesssim
r\iint_{\mathcal U_r}\sigma_+
\big(r^2|D_tE|+r|\nabla E|\big)\phi^2+o_r(1).
\]
The material-direction part is controlled by the transported-director equation:
\[
r^2|D_tE|\le C S_r+\mathcal R_{\mathrm{coh},\nu}^{E}.
\]
It is therefore bounded by the strain row and the coherent residual row.  For
spatial twist, on \(\{|\omega|>0\}\),
\[
|\nabla E|\le C\frac{|\nabla\omega|}{|\omega|}.
\]
Since \(\sigma_+\le |S_{\mathrm{loc}}||\omega|^2\), Young's inequality gives
\[
r^2\sigma_+\frac{|\nabla\omega|}{|\omega|}
\le
\frac12 r^3\sigma_+|S_{\mathrm{loc}}|
+C r|\nabla\omega|^2.
\]
After integration, the first term is \(\mathfrak B_{\mathrm{strain}}\) and the
second is \(\mathcal E_\omega\).  All absorber rows are components of \(D\).
\end{proof}

\begin{theorem}[K4-final forward absorption]
\label{thm:K4}
For
\[
\mathcal F_{\mathrm{phys}}^{\mathrm{final}}
=\mathcal F_{\mathrm{phys}}+R_{\mathrm{osc}}^{\mathrm{fam}},
\]
there is
\begin{equation}
\label{eq:Cforw-final}
C_{\mathrm{forw}}^{\mathrm{final}}
:=C_{\mathrm{old}}+C_{\mathrm{osc}}
\end{equation}
such that
\begin{equation}
\label{eq:K4-final}
P(r)\mathcal F_{\mathrm{phys}}^{\mathrm{final}}(r)
\le
C_{\mathrm{forw}}^{\mathrm{final}}D(r)+o_r(1).
\end{equation}
Together with K3,
\begin{equation}
\label{eq:K4-forward}
P(r)
\le
\frac{C_{\mathrm{forw}}^{\mathrm{final}}}{c_{\mathrm{final}}(\eta)}D(r)
+o_r(1).
\end{equation}
\end{theorem}

\begin{proof}
Split \(P\mathcal F_{\mathrm{phys}}^{\mathrm{final}}\) into the old physical
term and \(P R_{\mathrm{osc}}^{\mathrm{fam}}\), then apply
Lemmas~\ref{lem:K4-old} and~\ref{lem:K4-osc}.  Inequality
\eqref{eq:K4-forward} follows from \eqref{eq:cfinal}.
\end{proof}

\clearpage
\section{K5B: High-Strain Transport Rigidity}
\label{sec:K5B}

K5B is the main new analytic wall.  Standard deformation-gradient kinematics
give forward bounds, such as shape growth controlled by accumulated strain.
The needed direction is reverse: source-weighted high strain must force shape,
oscillation/twist, renewal, separation, leakage, quotient redundancy, or a
legitimate \(P\)-term.
This is the main analytic estimate required for closure.  The gates below
record the required estimates and constant routes; they are the burden to prove
rather than a citation-closed theorem already discharged in this note.

\begin{lemma}[K5B-0 interface bridge]
\label{lem:K5B-interface}
Let \(W\ge0\) be any K5-admissible source-weighted integrand, including
\(W=S_r\mathbf 1_{H_\Lambda}\).  For promoted branches with
core/collar/tail decomposition,
\[
\mathcal T_j=\mathcal T_j^{\mathrm{core}}\cup
\mathcal T_j^{\mathrm{collar}}\cup\mathcal T_j^{\mathrm{tail}},
\]
the source-weighted integral splits as
\[
r\iint_{\mathcal U_r}\sigma_+W\phi^2
\le
\sum_j r\int_{E_j}\int_{I_j}
\sigma_+W(\eta_{s_j,t}^{(j)}(a),t)\phi^2\,dt\,da
+C_W(\mathcal B_\omega+\mathcal D_{\mathrm{ARR}}
+\mathcal E_\nu^{\mathrm{coh},+})
+o_r(1).
\]
\end{lemma}

\begin{proof}
On each core, incompressibility gives \(\det D\eta_{s_j,t}^{(j)}=1\), so the
Lagrangian change of variables is exact.  The complement of the core is the
collar, tail, and coherent-residual set; by definition these pieces are charged
to the boundary, ARR/tail, and coherent residual rows.  This bridge contains no
core coercivity; it only separates the interface leakage.
\end{proof}

\setcounter{requiredlemma}{3}
\begin{requiredlemma}[Gate K5B-A: Source-weighted core matrix rigidity]
\label{lem:K5B-path}
Let \(F\) solve \(\dot F=(\nabla u)(\eta(t),t)F\) on a branch core with
\(\det F=1\), and let \(E(t)=e(t)\otimes e(t)\) be the source director.  For a
material label \(a\), define the source-time probability measure
\[
d\mu_{j,a}(t)
:=
\frac{r\,\sigma_+(\eta_{s_j,t}^{(j)}(a),t)
\phi^2(\eta_{s_j,t}^{(j)}(a),t)\,dt}
{P_{j,a}(r)+\varepsilon}.
\]
Let
\[
A_{j,a}^H
:=
\int_{\{S_r>\Lambda\}}
S_r(\eta_{s_j,t}^{(j)}(a),t)\,d\mu_{j,a}(t).
\]
Then there is \(C_{\mathrm{rig}}\) such that
\begin{equation}
\label{eq:path-rigidity}
A_{j,a}^H
\le
C_{\mathrm{rig}}\big(
R_{\mathrm{shape}}^{(j,a)}
+R_{\mathrm{osc}}^{(j,a)}
+1\big)+o_r(1).
\end{equation}
\end{requiredlemma}

\begin{proof}[Required route]
Partition the high-strain time set into intervals on which the projective
source director varies by at most a fixed small angle.  On a low-variation
interval, large source-weighted dimensionless strain forces Cauchy--Green
condition-number production because \(\det F=1\).  Intervals not covered by
this alternative are charged to projective direction variation.  The finite
partition cost becomes the \(+1\) term; after branch averaging this is a
legitimate \(P(r)\) contribution.
\end{proof}

\begin{requiredlemma}[Gate K5B-B: Coercive branch action]
\label{lem:K5B-branch}
Let
\[
A_j^H(r,\Lambda)
:=
\frac{1}{P_j(r)}
r\iint_{\{S_r>\Lambda\}\cap\mathcal T_j}
\sigma_+S_r\phi^2\,dx\,dt .
\]
Then
\begin{equation}
\label{eq:branch-action}
A_j^H
\le
C_{\mathrm{act}}\big(
R_{\mathrm{shape}}^{(j)}
+R_{\mathrm{osc}}^{(j)}
+R_{\mathrm{renew}}^{(j)}
+R_{\mathrm{sep}}^{(j)}
+1\big)+o_r(1).
\end{equation}
\end{requiredlemma}

\begin{proof}[Required route]
Use Lemma~\ref{lem:K5B-interface} to split the branch into core and leakage.
Loss of coherent identity is renewal.  Persistent but physically separating
branches are separation and then final physical charge.  On the remaining
nonrenewing, nonseparating core, Lemma~\ref{lem:K5B-path} gives shape or
oscillation/twist.  The \(+1\) term becomes \(P_j(r)\) after multiplication.
\end{proof}

\begin{requiredlemma}[Gate K5B-C: High-strain matrix bound]
\label{lem:K5B-bound}
Let \(H_\Lambda=\{S_r>\Lambda\}\cap\mathcal U_r\).  Assume K4-final and
Lemma~\ref{lem:K5B-branch}.  If \(a_{\mathrm{leak}}\) is the coefficient with
which K5B-0 leakage rows enter \(D\), then
\begin{equation}
\label{eq:K5B-bound}
r\iint_{H_\Lambda}\sigma_+S_r\phi^2
\le
a_HD(r)+C_HP(r)+o_r(1),
\qquad
a_H:=a_{\mathrm{leak}}+C_{\mathrm{act}}C_{\mathrm{forw}}^{\mathrm{final}}.
\end{equation}
No arbitrary smallness of \(a_H\) is asserted.
\end{requiredlemma}

\begin{proof}[Required route]
Multiply \eqref{eq:branch-action} by \(P_j(r)\), sum in \(j\), and use that
shape, oscillation, renewal, and separation are components of
\(\mathcal F_{\mathrm{phys}}^{\mathrm{final}}\).  K4-final routes their product
with \(P(r)\) into \(D(r)\).  Leakage contributes \(a_{\mathrm{leak}}D(r)\),
and the summed \(+1\) terms contribute \(C_HP(r)\).
\end{proof}

\begin{lemma}[Matrix-ready strain row]
\label{lem:strain-row}
The strain row satisfies
\begin{equation}
\label{eq:strain-linear}
\mathfrak B_{\mathrm{strain}}(r)
\le
a_{\mathrm{strain}}D(r)+C_{\mathrm{strain}}P(r)+o_r(1),
\qquad
a_{\mathrm{strain}}:=a_H,\quad
C_{\mathrm{strain}}:=\Lambda+C_H.
\end{equation}
\end{lemma}

\begin{proof}
Split \(\mathfrak B_{\mathrm{strain}}\) into \(\{S_r\le\Lambda\}\) and
\(H_\Lambda\).  The low-strain part is at most \(\Lambda P(r)\); the high-strain
part is Lemma~\ref{lem:K5B-bound}.
\end{proof}

\begin{remark}[Forbidden inference]
The kinematic estimate
\[
\frac{d}{dt}\log\kappa_j(t)\le C|S_{\mathrm{loc}}|
\]
does not imply
\[
\int |S_{\mathrm{loc}}|\,dt\lesssim \log\kappa_j .
\]
Large accumulated strain must instead be routed to shape, oscillation/twist,
renewal, separation, leakage, quotient redundancy, or \(P\).  Without the
oscillation/twist channel, K5B is incomplete.
\end{remark}

\clearpage
\section{K5: Reverse Matrix Absorption}
\label{sec:K5}

K5 packages all reverse row estimates into a nonnegative matrix inequality.
Use the row order in \eqref{eq:X-def}.  The required row form is:
\begin{align}
\mathcal E_\omega
&\le
a_{12}\mathcal B_\omega
+a_{13}\mathcal D_{\mathrm{ARR}}
+a_{14}\mathfrak B_{\mathrm{strain}}
+a_{16}\mathcal E_\nu^{\mathrm{coh},+}
+b_1P+o_r(1), \label{eq:row1}\\
\mathcal B_\omega
&\le
a_{21}\mathcal E_\omega
+a_{23}\mathcal D_{\mathrm{ARR}}
+a_{26}\mathcal E_\nu^{\mathrm{coh},+}
+b_2P+o_r(1), \label{eq:row2}\\
\mathcal D_{\mathrm{ARR}}
&\le
a_{31}\mathcal E_\omega
+a_{32}\mathcal B_\omega
+a_{36}\mathcal E_\nu^{\mathrm{coh},+}
+b_3P+o_r(1), \label{eq:row3}\\
\mathfrak B_{\mathrm{strain}}
&\le
a_{\mathrm{strain}}\mathbf 1^TX
+C_{\mathrm{strain}}P+o_r(1), \label{eq:row4}\\
M_{\mathrm{deact}}
&\le
a_{51}\mathcal E_\omega
+a_{52}\mathcal B_\omega
+a_{53}\mathcal D_{\mathrm{ARR}}
+a_{54}\mathfrak B_{\mathrm{strain}}
+a_{55}M_{\mathrm{deact}}
+a_{56}\mathcal E_\nu^{\mathrm{coh},+}
+b_5P+o_r(1), \label{eq:row5}\\
\mathcal E_\nu^{\mathrm{coh},+}
&\le
a_{61}\mathcal E_\omega
+a_{62}\mathcal B_\omega
+a_{63}\mathcal D_{\mathrm{ARR}}
+b_6P+o_r(1). \label{eq:row6}
\end{align}

Equivalently,
\begin{equation}
\label{eq:AXb}
X(r)\le A X(r)+bP(r)+o_r(1),
\end{equation}
where
\[
A=
\begin{pmatrix}
0&a_{12}&a_{13}&a_{14}&0&a_{16}\\
a_{21}&0&a_{23}&0&0&a_{26}\\
a_{31}&a_{32}&0&0&0&a_{36}\\
a_{\mathrm{strain}}&a_{\mathrm{strain}}&a_{\mathrm{strain}}
&a_{\mathrm{strain}}&a_{\mathrm{strain}}&a_{\mathrm{strain}}\\
a_{51}&a_{52}&a_{53}&a_{54}&a_{55}&a_{56}\\
a_{61}&a_{62}&a_{63}&0&0&0
\end{pmatrix},
\qquad
b=
\begin{pmatrix}
b_1\\ b_2\\ b_3\\ C_{\mathrm{strain}}\\ b_5\\ b_6
\end{pmatrix}.
\]

\begin{theorem}[Matrix-ready reverse absorption]
\label{thm:K5}
Assume the row estimates \eqref{eq:row1}--\eqref{eq:row6} and
\[
\rho(A)<1.
\]
Then
\begin{equation}
\label{eq:reverse}
D(r)\le C_0P(r)+o_r(1),
\qquad
C_0:=\mathbf 1^T(I-A)^{-1}b.
\end{equation}
\end{theorem}

\begin{proof}
Since \(A\ge0\) and \(\rho(A)<1\), the Neumann series
\((I-A)^{-1}=\sum_{k\ge0}A^k\) converges and is nonnegative.  Applying it
componentwise to \eqref{eq:AXb} gives
\[
X(r)\le (I-A)^{-1}bP(r)+o_r(1).
\]
Multiplying by \(\mathbf 1^T\) gives \eqref{eq:reverse}.
\end{proof}

For certification without explicitly inverting \(I-A\), let \(w_i>0\) and set
\[
q_A(w):=\max_i\frac{1}{w_i}\sum_j A_{ij}w_j.
\]
Equivalently, the six weighted row sums are the row fractions of \(A\) in the
weighted sup norm.  Define
\begin{equation}
\label{eq:C0-cert}
C_0^{\mathrm{cert}}(w)
:=
\frac{
(\sum_iw_i)\max_i(b_i/w_i)
}{
1-q_A(w)
}.
\end{equation}

\begin{lemma}[Weighted row-sum certificate]
\label{lem:weighted-certificate}
If \(q_A(w)<1\), then \(\rho(A)\le q_A(w)<1\), and
\[
\mathbf 1^T(I-A)^{-1}b\le C_0^{\mathrm{cert}}(w).
\]
\end{lemma}

\begin{proof}
Use the weighted sup norm \(\|y\|_w=\max_i |y_i|/w_i\).  The operator norm of
\(A\) in this norm is \(q_A(w)\), so \(\rho(A)\le q_A(w)\).  If
\(\beta_w:=\max_i(b_i/w_i)\), then \((A^kb)_i\le \beta_w q_A(w)^kw_i\).
Summing over rows and \(k\ge0\) gives \eqref{eq:C0-cert}.
\end{proof}

\begin{remark}[Coefficient audit]
Every coefficient in \(A\) must come from a row proof.  In particular, the
strain row contributes
\[
A_{4j}=a_{\mathrm{strain}}\quad(1\le j\le6),
\]
so K5B constants directly affect the feasibility of \(q_A(w)<1\).  K5 does not
permit unproved smallness assumptions.
\end{remark}

\clearpage
\section{K6: Spectral and Loop-Gap Certification}
\label{sec:K6}

K6 contains no PDE estimate.  It is the strict algebraic loop that combines K3,
K4, and K5.

\begin{theorem}[Absorption loop algebra]
\label{thm:K6}
Assume
\begin{equation}
\label{eq:loop-forward}
P(r)\le
\frac{C_{\mathrm{forw}}^{\mathrm{final}}}{c_{\mathrm{final}}(\eta)}
D(r)+o_r(1),
\end{equation}
\begin{equation}
\label{eq:loop-reverse}
D(r)\le C_0P(r)+o_r(1),
\end{equation}
and the strict gap
\begin{equation}
\label{eq:loop-gap}
C_{\mathrm{forw}}^{\mathrm{final}}C_0<c_{\mathrm{final}}(\eta).
\end{equation}
Then along the contradiction sequence,
\[
P(r_n)\to0,\qquad D(r_n)\to0.
\]
\end{theorem}

\begin{proof}
Let
\[
q:=\frac{C_{\mathrm{forw}}^{\mathrm{final}}C_0}
{c_{\mathrm{final}}(\eta)}<1.
\]
Combining \eqref{eq:loop-forward} and \eqref{eq:loop-reverse} gives
\[
D(r_n)\le qD(r_n)+o_n(1),
\]
and hence \(D(r_n)\to0\).  Then \eqref{eq:loop-forward} gives
\(P(r_n)\to0\).
\end{proof}

\begin{corollary}[Certified loop-gap condition]
\label{cor:cert-loop}
If \(q_A(w)<1\) and
\begin{equation}
\label{eq:cert-loop-gap}
C_{\mathrm{forw}}^{\mathrm{final}}C_0^{\mathrm{cert}}(w)
<
c_{\mathrm{final}}(\eta),
\end{equation}
then the K6 loop closes with the certified reverse constant
\(C_0^{\mathrm{cert}}(w)\).
\end{corollary}

\begin{proof}
Lemma~\ref{lem:weighted-certificate} supplies
\[
D(r)\le C_0^{\mathrm{cert}}(w)P(r)+o_r(1).
\]
Condition \eqref{eq:cert-loop-gap} is exactly the strict K6 gap using this
certified reverse constant.
\end{proof}

The single constant inequality to verify is therefore
\begin{equation}
\label{eq:single-ledger}
C_{\mathrm{forw}}^{\mathrm{final}}\mathbf 1^T(I-A)^{-1}b
<
\frac{\delta c_m^{\mathrm{final}}}{C_*^{\mathrm{final}}}\eta.
\end{equation}
The checkable sufficient version is
\[
q_A(w)<1,
\qquad
C_{\mathrm{forw}}^{\mathrm{final}}C_0^{\mathrm{cert}}(w)
<
c_{\mathrm{final}}(\eta).
\]
No qualitative positivity statement can replace this strict gap.

\clearpage
\section{K7: Compatibility and Regularity Bridge}
\label{sec:K7}

K7 is standard regularity theory in the exact denominator normalization used by
the ledger.  Its burden is compatibility: every edit to \(D\) must still imply
smallness of the packages used below.

\begin{definition}[K7 local packages]
\label{def:K7-packages}
Let \(\phi\in C_c^\infty(B_{2r})\) satisfy \(\phi\equiv1\) on \(B_r\), and set
\[
A_\omega(2r)
:=
2r\,\operatorname*{ess\,sup}_{t\in(T-(2r)^2,T)}
\int_{B_{2r}}|\omega|^2\phi^2\,dx,
\]
\[
B_\omega^{\mathrm{diss}}(2r)
:=
\nu\,2r\iint_{Q_{2r}}|\nabla\omega|^2\phi^2\,dx\,dt,
\]
\[
K_{\omega,3}^{\mathrm{crit}}(r;z_0)
:=
r\iint_{Q_r(z_0)}|\omega|^3\,dx\,dt.
\]
The compatibility condition is
\begin{equation}
\label{eq:K7-compat}
A_\omega(2r)+B_\omega^{\mathrm{diss}}(2r)
\le CD(r)+o_r(1),
\qquad
\mathcal H_u(r)+\mathcal H_p(r)+\mathcal R_\phi(r)
\le CD(r)+o_r(1),
\end{equation}
where \(\mathcal H_u,\mathcal H_p\) are local harmonic/projection remainders
and \(\mathcal R_\phi\) is the cutoff remainder.
\end{definition}

\begin{lemma}[Vorticity interpolation]
\label{lem:vort-interp}
There is \(C\) such that
\begin{equation}
\label{eq:vort-interp}
K_{\omega,3}^{\mathrm{crit}}(r;z_0)
\le
C\nu^{-3/4}
A_\omega(2r)^{3/4}B_\omega^{\mathrm{diss}}(2r)^{3/4}
+\mathcal R_\phi(r).
\end{equation}
\end{lemma}

\begin{proof}
This is the scale-correct local Gagliardo--Nirenberg estimate applied to
\(\omega\phi\), integrated over a parabolic interval and multiplied by \(r\).
\end{proof}

\begin{lemma}[Mean-subtracted velocity and pressure control]
\label{lem:velocity-pressure}
The standard local div--curl/Biot--Savart and pressure decomposition estimates
give
\[
r^{-2}\iint_{Q_r}|u-\bar u_{B_r}(t)|^3
\le
C K_{\omega,3}^{\mathrm{crit}}(2r;z_0)+\mathcal H_u(r),
\]
and
\[
r^{-2}\iint_{Q_r}|p-\bar p_{B_r}(t)|^{3/2}
\le
C r^{-2}\iint_{Q_{2r}}|u-\bar u_{B_{2r}}(t)|^3+\mathcal H_p(r).
\]
\end{lemma}

\begin{proof}
Subtracting the ball mean fixes the local Galilean gauge.  The localized
Biot--Savart piece is controlled by the \(L^3\) vorticity norm; the harmonic
velocity remainder is \(\mathcal H_u\).  For pressure, write
\(-\Delta p_{\mathrm{loc}}
=\partial_i\partial_j[(u_i-\bar u_i)(u_j-\bar u_j)]\) on a localized ball and
apply Calderon--Zygmund estimates.  The harmonic pressure remainder is
\(\mathcal H_p\).
\end{proof}

\begin{lemma}[Mean-subtracted CKN epsilon criterion]
\label{lem:CKN}
There exists \(\varepsilon_{\mathrm{CKN}}>0\) such that
\[
\mathcal K_{\mathrm{CKN}}^{\mathrm{ms}}(r;z_0)<\varepsilon_{\mathrm{CKN}}
\]
implies that \(z_0\) is regular.
\end{lemma}

\begin{theorem}[K7 regularity bridge]
\label{thm:K7}
Assume \eqref{eq:K7-compat}.  If \(D(r_n)\to0\) along \(r_n\downarrow0\), then
\[
\mathcal K_{\mathrm{CKN}}^{\mathrm{ms}}(r_n;z_0)\to0.
\]
Consequently \(z_0\) is regular.
\end{theorem}

\begin{proof}
Compatibility makes the localized vorticity energy-dissipation package and all
harmless remainders vanish.  Lemma~\ref{lem:vort-interp} gives
\(K_{\omega,3}^{\mathrm{crit}}(2r_n;z_0)\to0\).  Lemma~\ref{lem:velocity-pressure}
then gives mean-subtracted velocity and pressure smallness, so
\(\mathcal K_{\mathrm{CKN}}^{\mathrm{ms}}(r_n;z_0)\to0\).  Apply
Lemma~\ref{lem:CKN}.
\end{proof}

\clearpage
\section{Contradiction Chain}
\label{sec:chain}

\begin{theorem}[Conditional K3--K7 closure]
\label{thm:closure}
Assume the upstream promoted-source package, K3, K4, K5/K5B, K6, and K7 in the
forms stated above.  If the strict loop gap \eqref{eq:loop-gap} holds, then the
candidate point \(z_0\) is regular.
\end{theorem}

\begin{proof}
If the promoted family is quotient-redundant, it is not a nonredundant
amplification channel.  Otherwise K3 gives
\[
\mathcal F_{\mathrm{phys}}^{\mathrm{final}}\ge c_{\mathrm{final}}(\eta).
\]
K4 gives
\[
P\le
\frac{C_{\mathrm{forw}}^{\mathrm{final}}}{c_{\mathrm{final}}(\eta)}D+o.
\]
K5 gives \(D\le C_0P+o\), with either
\[
C_0=\mathbf 1^T(I-A)^{-1}b
\]
or a certified sufficient constant \(C_0^{\mathrm{cert}}(w)\).  K6 and the
strict gap force \(P(r_n)\to0\) and \(D(r_n)\to0\).  K7 then converts
\(D(r_n)\to0\) into CKN smallness and regularity, contradicting singularity.
\end{proof}

\clearpage
\section{Falsification and Attack Surface}
\label{sec:attack-surface}

This note is useful only if it makes the failure modes visible.  A critic can
break the K3--K7 closure by showing any one of the following:
\begin{enumerate}
\item A nonredundant promoted source sequence with positive \(P(r_n)\) and no
final physical payment.
\item High source-weighted strain that avoids shape, oscillation/twist,
renewal, separation, leakage, quotient redundancy, and a legitimate \(P\)-term.
\item A K5 matrix whose proved coefficients cannot satisfy
\(\rho(A)<1\), or cannot satisfy a certified weighted row condition
\(q_A(w)<1\).
\item A failure of the strict K6 loop gap
\[
C_{\mathrm{forw}}^{\mathrm{final}}C_0<c_{\mathrm{final}}(\eta).
\]
\item A denominator \(D\) that collapses without implying mean-subtracted CKN
smallness.
\end{enumerate}
These are the attack points; if one survives, the proposed closure survives
only as a conditional ledger, not as a completed regularity proof.

\clearpage
\section{Proof-Burden Ledger}
\label{sec:burden}

\begin{longtable}{@{}p{0.12\linewidth}p{0.34\linewidth}p{0.42\linewidth}@{}}
\toprule
\textbf{Gate} & \textbf{Claim} & \textbf{Remaining burden} \\
\midrule
\endhead
K3 &
Nonredundant promoted packing either quotient-collapses or yields
\(\mathcal F_{\mathrm{phys}}^{\mathrm{final}}\ge c_{\mathrm{final}}(\eta)\). &
Graph quotient constants, finite-overlap mass transfer, final physical
coordinate thresholds, and the oscillation-adjoint estimate. \\
\addlinespace
K4 &
\(P\mathcal F_{\mathrm{phys}}^{\mathrm{final}}\le
C_{\mathrm{forw}}^{\mathrm{final}}D+o\). &
Forward routing of all physical ledgers, especially \(R_{\mathrm{osc}}^{\mathrm{fam}}\),
into the chosen denominator rows with scale-independent constants. \\
\addlinespace
K5B &
Source-weighted high strain routes to shape, oscillation, renewal, separation,
leakage, quotient redundancy, or \(P\). &
The reverse high-strain rigidity estimate
\eqref{eq:K5B-bound}; no substitution of an unweighted path diagnostic for the
source-weighted action. \\
\addlinespace
K5 &
The six denominator rows satisfy \(X\le AX+bP+o\), and \(\rho(A)<1\). &
Row-by-row coefficient provenance, packed boundary constants, ARR/tail routing,
deactivation and coherent-residual constants, and a spectral certificate such
as \(q_A(w)<1\). \\
\addlinespace
K6 &
The constants satisfy
\(C_{\mathrm{forw}}^{\mathrm{final}}C_0<c_{\mathrm{final}}(\eta)\). &
Strict numerical loop gap; positivity of \(c_{\mathrm{final}}(\eta)\) alone is
not enough. \\
\addlinespace
K7 &
\(D(r_n)\to0\) implies mean-subtracted CKN smallness. &
Compatibility of the exact denominator rows with local vorticity
energy-dissipation, harmless cutoff/projection remainders, and the
mean-subtracted CKN normalization. \\
\bottomrule
\end{longtable}

\section{Conclusion}
\label{sec:conclusion}

This trimmed note is the public technical burden behind the imported gates in
the amplification-payment paper.  The proof is not reduced to one inequality
until every row coefficient in K5 is proved, the K5B high-strain transport wall
is crossed, and K6 verifies the strict constant gap.  K7 is standard in spirit,
but it remains coupled to the exact denominator normalization: edits to the
ledger must preserve \eqref{eq:K7-compat}.

\end{document}
